\chapter{结论与展望}

\section{研究工作总结}
本文围绕动态数据环境下概念漂移导致的哈希检索性能退化问题，分别在单模态与跨模态场景下开展了系统研究，形成了由概念保留到多表增量学习的连续研究思路。针对传统静态哈希方法在类别扩展与分布演化条件下难以保持长期稳定性能的问题，本文重点关注了动态场景中的语义保留、模型增量更新与检索效率之间的协同关系。

在单模态图像检索方面，本文提出了半监督概念保留哈希方法（SCPH）。该方法通过固定类别概念编码，将新到达样本持续投影到稳定语义区域，并结合 KNN 伪标签机制利用未标注样本结构信息，在标注受限条件下提升模型可学习性。与此同时，SCPH在模型更新过程中避免对历史数据哈希码进行重编码，兼顾了动态环境下的检索性能与更新效率。围绕新类出现与数据分布改变两类概念漂移情形的实验结果表明，SCPH在不同标注比例下均表现出较好的稳定性与竞争性能。

在跨模态检索方面，本文提出了基于多表结构的增量式跨模态哈希方法（MIH）。该方法通过构建固定规模的多哈希表系统，在时间维度上并行保留不同阶段学习到的跨模态语义结构，缓解单模型持续更新带来的历史知识覆盖问题。进一步地，MIH设计了基于新到达数据表现的哈希表权重更新机制，并引入查询自适应加权策略，对检索阶段的汉明距离进行细粒度调节，从而提升动态跨模态场景下的语义判别能力。实验与消融结果验证了多表机制及两类加权策略对性能提升的有效性。

在系统应用方面，本文第五章面向电商场景完成了基于在线哈希的智能检索系统设计与实现。该系统以 SCPH 与 MIH 为核心引擎，采用数据接入层、特征工程层、核心哈希引擎层、索引存储层与业务服务层的五层架构，实现了相似商品推荐、拍照购与自然语言搜索三类核心功能，打通了“数据接入—特征提取—哈希编码—索引检索—业务服务”的端到端链路。通过系统化测试验证了 SCPH 与 MIH 在工程环境中的可落地性，以及在不重编码历史哈希码前提下的在线增量更新与检索稳定性，将前述算法研究成果转化为可部署的系统方案，为在线哈希方法在产业级检索系统中的应用提供了实践依据。

总体而言，本文在概念漂移条件下在线哈希检索问题上实现了从单模态到跨模态、从算法方法到系统应用的完整研究拓展，建立了面向非平稳数据环境的在线哈希检索方法体系，为后续研究提供了可复用的研究思路与实验范式。

\section{主要研究结论}
综合第三、四、五章的研究与实验结果可以得出，概念漂移并非仅表现为数据分布的局部波动，而是会在时间维度上持续扰动哈希空间中的语义邻域结构，进而引发检索性能的系统性退化。本文的核心结论是：在动态检索任务中，只有同时处理“跨时间语义保持”与“增量更新适配”这两个目标，才能在长期运行条件下维持稳定有效的哈希检索性能。围绕这一认识，本文在单模态与跨模态两个层面分别给出了可验证的实现路径，并在对应实验场景中获得了一致支持；进一步地，第五章所实现的电商智能检索系统表明，上述方法能够嵌入真实业务链路，在“数据接入—特征提取—哈希编码—索引检索—业务服务”全流程中保持免重编码的在线更新与稳定检索，验证了从算法到工程落地的可行性。

在单模态动态检索场景下，SCPH 的研究表明，基于固定概念编码的概念保留约束能够显著增强模型对概念漂移的抗扰动能力。通过将类别语义映射为汉明空间中的稳定参照，并在后续时间步持续约束新样本向目标语义区域对齐，SCPH在不重编码历史样本的前提下实现了检索性能与更新效率的兼顾。进一步地，引入基于邻域关系的伪标签推断后，未标注样本能够为模型提供有效的结构信息补充，使方法在标注比例受限时仍保持较好的鲁棒性。基于 CIFAR-10与 CIFAR-100构建的新类出现与分布变化两类漂移场景的结果显示，SCPH在动态演化过程中能够更稳定地保持检索精度，验证了“概念保留+半监督”在单模态动态哈希中的有效性。

在跨模态动态检索场景下，MIH 的研究进一步说明，单一模型持续更新难以同时保持历史语义与适配当前分布，而多表增量学习能够更好地缓解这一结构性矛盾。MIH通过并行维护不同时段学习得到的跨模态哈希表，避免新知识学习对历史语义映射的过度覆盖，从机制上降低了灾难性遗忘风险；同时，通过基于最新数据表现的表权重更新与查询自适应加权，模型能够在检索阶段动态提升更具判别力的表与比特位贡献，从而提高排序质量。MIRFlickr 与 MSCOCO 两类概念漂移场景下的双向检索结果及消融实验共同表明，多表保留机制与双重加权策略对提升跨模态动态检索稳定性具有关键作用。

总体而言，本文形成了由单模态到跨模态、由概念保留到多表增量的连续研究结论：概念漂移条件下动态哈希检索的性能提升，不应依赖单一强更新，而应建立在语义锚定、历史信息保留与动态判别增强的协同机制之上。该结论不仅解释了现有方法在长期动态环境中易退化的原因，也为后续研究面向动态数据的持续学习型检索模型提供了明确方向。

\section{不足与展望}
尽管本文提出的SCPH与MIH在新类出现与数据分布变化两类概念漂移场景下均取得了较为稳定的检索性能，但从动态检索系统的长期运行需求出发，现阶段方法仍存在进一步完善空间。首先，在更新机制方面，本文两类方法均采用按时间批次触发模型更新的策略，尚未将“是否发生显著漂移”作为更新触发条件显式纳入框架。该设定有助于实验对比，但在真实在线环境中可能引入不必要更新。后续可将概念漂移检测与增量学习过程耦合，构建自适应更新机制，以在保证性能的同时进一步降低计算与维护成本。

其次，在监督信息利用方面，SCPH通过邻域伪标签缓解了标注稀缺问题，但伪标签质量仍受样本分布复杂度与类别边界清晰度影响。在分布重叠较强或噪声较大的场景中，误标传播可能削弱模型稳定性。未来可进一步引入置信度评估、样本不确定性建模与一致性约束机制，提升伪标签的可用性与训练过程的稳健性，从而增强方法在弱监督动态环境中的泛化能力。

另外，在跨模态增量学习方面，MIH通过固定规模多表机制实现了历史信息保留与当前分布适配之间的平衡，但在更大规模、长周期场景下，表规模设定与权重更新策略仍存在优化空间。后续研究可围绕自适应表规模控制、跨表知识蒸馏与轻量化检索计算展开，以在不牺牲性能的前提下提升系统可扩展性。此外，本文目前主要基于模态同步到达设定，尚未覆盖真实场景中更常见的模态异步到达、模态缺失及弱配对监督条件。将MIH扩展至上述复杂设定，将是提升方法实用价值的重要方向。

最后，在理论层面，本文已通过实验验证了方法有效性，但对动态环境下的收敛行为、误差传播路径与性能边界仍缺乏更系统的理论分析与研究。未来可进一步从优化理论与统计学习角度建立分析框架，增强模型在概念漂移条件下的可解释性与理论完备性。在系统应用层面，第五章所实现的电商智能检索系统已在端到端链路上验证了 SCPH 与 MIH 的可用性，但当前规模与负载仍属演示与功能验证阶段；后续可在更大商品规模与更高并发下开展压力测试与性能优化，推动方法向产业级部署进一步迈进。总体而言，本文为概念漂移下的在线哈希检索提供了可行的研究基础，后续工作可围绕自适应更新、鲁棒学习、异步多模态建模与理论深化持续推进，以支持复杂真实环境中的长期稳定检索应用。
